% !TEX root = ../main.tex
\corrigesChapitre{4}{Variations et optimisation}
\corrigesSouples

\begin{blocLivre}{cor-4-1}
\begin{corrige}{1}{Du signe aux variations}
\demarche{Traduire les signes de la dérivée en flèches et comparer aussi les valeurs aux bornes.}
\begin{questions}\item La fonction croît sur $[-3;1]$, puis décroît sur $[1;4]$.
Son maximum est atteint en 1 et vaut $\res{7}$.
\item On compare les deux bornes : $-1<2$, donc le minimum global vaut
$\res{-1}$, atteint en 4.
\end{questions}
\label{figlivre-r04-01-var-cor}
\figureLivre[corrige]{r04-01-var}

\end{corrige}
\end{blocLivre}
\begin{blocLivre}{cor-4-2}
\begin{corrige}{2}{Un trinôme sur un segment}
\demarche{Chercher le zéro de la dérivée à l’intérieur du segment.}
$f'(x)=2x-4$, négative pour $x<2$, nulle en 2, positive pour $x>2$.
Donc $f$ décroît sur $[-1;2]$, de $f(-1)=6$ à $f(2)=-3$, puis croît
sur $[2;5]$ jusqu'à $f(5)=6$.
Le minimum est $\res{-3}$ en 2 ; le maximum est $\res{6}$, aux deux bornes $-1$ et 5.
\label{figlivre-r04-02-var-cor}
\figureLivre[corrige]{r04-02-var}

\end{corrige}
\end{blocLivre}
\begin{blocLivre}{cor-4-3}
\begin{corrige}{3}{Une dérivée ne change pas de signe}
\demarche{Une dérivée strictement positive impose une croissance stricte.}
$g'(x)=3x^2+3>0$ pour tout réel $x$, donc $g$ est strictement croissante sur $\R$.
Sur $[-1;2]$, le minimum vaut $g(-1)=\res{-4}$ et le maximum $g(2)=\res{14}$.
\label{figlivre-r04-03-var-cor}
\figureLivre[corrige]{r04-03-var}

\end{corrige}
\end{blocLivre}
\begin{blocLivre}{cor-4-4}
\begin{corrige}{4}{Un quotient simple}
\demarche{Le dénominateur de la dérivée est positif sur tout le domaine étudié.}
$h'(x)=\dfrac{x+2-(x+1)}{(x+2)^2}=\dfrac1{(x+2)^2}>0$ sur $[0;4]$.
La fonction croît de $\res{h(0)=1/2}$ à $\res{h(4)=5/6}$ : ce sont son
minimum et son maximum.
\label{figlivre-r04-04-var-cor}
\figureLivre[corrige]{r04-04-var}

\end{corrige}
\end{blocLivre}
\begin{blocLivre}{cor-4-5}
\begin{corrige}{5}{Dérivée nulle et extremum}
\demarche{Un zéro isolé de la dérivée n’entraîne un extremum que si le sens de variation change.}
$f'(x)=3x^2\geq0$ et $f'(0)=0$. La dérivée est positive sauf en 0 :
$f$ est strictement croissante sur $\R$. Elle ne change pas de sens de variation
en 0 : $\res{\text{aucun extremum en }0}$.
On peut aussi comparer $f(-\varepsilon)<f(0)<f(\varepsilon)$ pour tout $\varepsilon>0$.
\label{figlivre-r04-05-var-cor}
\figureLivre[corrige]{r04-05-var}

\end{corrige}
\end{blocLivre}
\begin{blocLivre}{cor-4-6}
\begin{corrige}{6}{Un polynôme du troisième degré}
\demarche{Factoriser la dérivée et distinguer les valeurs locales des valeurs globales.}
\begin{questions}
\item $f'(x)=3x^2-6x-9=3(x^2-2x-3)=\res{3(x+1)(x-3)}$.
\item Le produit est positif à l'extérieur des racines et négatif entre elles.
\end{questions}
\begin{center}\begin{tableau}
\otTab[2.5]{$x$/0.8,$f'(x)$/0.8,$f$/1.8}{$-3$,$-1$,$3$,$5$}
\tkzTabLine{,+,z,-,z,+,}
\tkzTabVar{-/$-22$,+/$10$,-/$-22$,+/$10$}
\end{tableau}\end{center}
\otLabel{3.}\enspace À l'intérieur du segment, maximum local 10 en $-1$,
minimum local $-22$ en 3. Le maximum global 10 est atteint en $-1$ \emph{et} 5 ;
le minimum global $-22$ est atteint en $-3$ \emph{et} 3.
\end{corrige}
\end{blocLivre}
\begin{blocLivre}{cor-4-7}
\begin{corrige}{7}{Un minimum avec un quotient}
\demarche{Sur les réels positifs, isoler le seul facteur qui peut changer de signe.}
\begin{questions}
\item $f'(x)=1-4/x^2=\dfrac{(x-2)(x+2)}{x^2}$.
Sur $]0;+\infty[$, son signe est celui de $x-2$.
\item La fonction décroît jusqu'à 2, puis croît. Son minimum est
$\res{f(2)=4}$ ; donc $x+4/x\geq4$ pour tout $x>0$.
\item $f(x)-4=(x^2-4x+4)/x=\res{(x-2)^2/x}\geq0$ car $x>0$.
L'égalité a lieu exactement pour $x=2$.
\end{questions}
\label{figlivre-r04-07-var-cor}
\figureLivre[corrige]{r04-07-var}

\end{corrige}
\end{blocLivre}
\begin{blocLivre}{cor-4-8}
\begin{corrige}{8}{Étudier une fonction paire}
\demarche{Étudier une moitié du domaine puis utiliser la symétrie.}
\begin{questions}
\item $f(-x)=(-x)^4-4(-x)^2=f(x)$ : $f$ est paire.
\item $f'(x)=4x(x^2-2)$. Sur $[0;3]$, elle est négative entre 0 et $\sqrt2$,
puis positive après $\sqrt2$ : $f$ décroît de 0 à $-4$, puis croît jusqu'à 45.
\item Sur $[-3;0]$, par symétrie : décroissance de 45 à $-4$ entre $-3$ et $-\sqrt2$,
puis croissance de $-4$ à 0. Minimum global $\res{-4}$ en $\pm\sqrt2$ ; maximum
global $\res{45}$ en $\pm3$.
\item $g(-x)=-x^3+3x=-g(x)$ : $g$ est impaire et sa courbe est symétrique par rapport à l'origine.
\end{questions}
\label{figlivre-r04-08-var-cor}
\figureLivre[corrige]{r04-08-var}

\end{corrige}
\end{blocLivre}
\begin{blocLivre}{cor-4-9}
\begin{corrige}{9}{Position par rapport à une tangente}
\demarche{La différence entre la fonction et sa tangente donne leur position relative.}
\begin{questions}
\item $f(1)=1$, $f'(1)=3$ : $\res{T:y=3x-2}$.
\item $(x-1)^2(x+2)=(x^2-2x+1)(x+2)=x^3-3x+2$.
\item Sur $[0;3]$, $(x-1)^2\geq0$ et $x+2>0$, donc
$f(x)-(3x-2)\geq0$. La courbe est au-dessus de $T$ et la rencontre uniquement
au point $\res{(1;1)}$ sur cet intervalle.
\end{questions}
\label{figlivre-r04-09-position-cor}
\figureLivre[corrige]{r04-09-position}

\end{corrige}
\end{blocLivre}
\begin{blocLivre}{cor-4-10}
\begin{corrige}{10}{Une inégalité grâce aux variations}
\demarche{Mettre la dérivée au même dénominateur pour lire son signe.}
\begin{questions}
\item $f'(x)=1/(2\sqrt x)-1/2=\dfrac{1-\sqrt x}{2\sqrt x}$.
Pour $x>1$, le dénominateur est positif et le numérateur négatif.
Donc $f$ est strictement décroissante sur $[1;9]$.
\item $f(1)=1/2$ et $f(9)=3-9/2=-3/2$. Maximum $1/2$ en 1 ; minimum $-3/2$ en 9.
\item Toute valeur de $f$ est comprise entre ces extremums :
$\res{-\frac32\leq\sqrt x-\frac x2\leq\frac12}$.
\end{questions}
\label{figlivre-r04-10-var-cor}
\figureLivre[corrige]{r04-10-var}

\end{corrige}
\end{blocLivre}
\begin{blocLivre}{cor-4-11}
\begin{corrige}{11}{Optimiser une quantité entière}
\demarche{Chercher l’optimum réel, puis comparer les entiers qui l’encadrent.}
\begin{questions}
\item $B'(x)=-4x+27$ s'annule en $27/4=6{,}75$.
Le bénéfice croît jusqu'à 6,75 puis décroît.
\item Les deux entiers voisins sont $\res{6\text{ et }7}$.
Tous les autres entiers sont moins favorables d'après les variations.
\item $B(6)=50$ et $B(7)=51$. La quantité optimale est $\res{7\text{ articles}}$,
pour un bénéfice de $\res{51\text{ euros}}$.
\end{questions}
\label{figlivre-r04-11-var-cor}
\figureLivre[corrige]{r04-11-var}

\end{corrige}
\end{blocLivre}
\begin{blocLivre}{cor-4-12}
\begin{corrige}{12}{Une boîte sans couvercle}
\demarche{Exprimer le volume avec une seule variable et contrôler les dimensions positives.}
\begin{questions}
\item La hauteur est $x>0$ et le fond a pour côté $24-2x>0$ : $0<x<12$.
Le volume vaut $\res{V(x)=x(24-2x)^2}$ cm$^3$.
\item En développant, $V(x)=4x^3-96x^2+576x$.
Donc $V'(x)=12x^2-192x+576=\res{12(x-4)(x-12)}$.
La dérivée est positive sur $]0;4[$ et négative sur $]4;12[$.
Le volume croît jusqu'à 4 puis décroît.
\item Le maximum est atteint pour $\res{x=4\text{ cm}}$.
Le fond mesure $16\times16$ cm et la hauteur 4 cm.
Le volume est $4\times16^2=1024$ cm$^3$, soit $\res{1{,}024\text{ L}}$.
\end{questions}
\label{figlivre-r04-12-var-cor}
\figureLivre[corrige]{r04-12-var}

\label{figlivre-r04-12-boite-cor}
\figureLivre[corrige]{r04-12-boite}

\end{corrige}
\end{blocLivre}
\begin{blocLivre}{cor-4-13}
\begin{corrige}{13}{Un rectangle sous une parabole}
\demarche{Relier les coordonnées des sommets à la largeur et à la hauteur.}
\begin{questions}
\item La largeur vaut $2x$ et la hauteur $9-x^2>0$, donc
$\res{A(x)=2x(9-x^2)=18x-2x^3}$.
\item $A'(x)=18-6x^2=6(\sqrt3-x)(\sqrt3+x)$.
Sur $]0;3[$, le signe est positif avant $\sqrt3$, négatif après.
Le maximum est donc atteint en $\res{x=\sqrt3}$.
\item La largeur est $\res{2\sqrt3}$, la hauteur $\res{6}$ et l'aire
$\res{12\sqrt3\approx20{,}78}$ unités d'aire.
\end{questions}
\begin{center}\begin{tikzpicture}\begin{axis}[otLivreGraphe,width=91mm,height=48mm,xmin=-3.3,xmax=3.3,ymin=-0.5,ymax=10,xtick={-3,0,3},ytick={3,6,9}]
\addplot[otNavy,thick,domain=-3:3,samples=70]{9-x^2};
\draw[otGreenDark,fill=otGreen,fill opacity=.13](axis cs:-1.73205,0)rectangle(axis cs:1.73205,6);
\end{axis}\end{tikzpicture}\end{center}
\label{figlivre-r04-13-var-cor}
\figureLivre[corrige]{r04-13-var}

\end{corrige}
\end{blocLivre}
\begin{blocLivre}{cor-4-14}
\begin{corrige}{14}{Un emballage de périmètre imposé}
\demarche{Utiliser le périmètre pour éliminer la hauteur avant de dériver le volume.}
\begin{questions}
\item $2r+2h=40$, donc $h=20-r$. Les deux dimensions étant strictement positives,
$\res{0<r<20}$.
\item $V(r)=\pi r^2(20-r)$, donc $\res{V'(r)=\pi r(40-3r)}$.
\item Sur le domaine, le signe dépend de $40-3r$ : $V$ croît jusqu'à $40/3$,
puis décroît. Le rayon optimal est $\res{40/3\text{ cm}}$ et la hauteur $\res{20/3\text{ cm}}$.
\item Le volume maximal vaut $\res{32000\pi/27\text{ cm}^3}$, environ 3723,37 cm$^3$.
Aux cas limites, le rayon ou la hauteur est nul : le volume est nul et l'objet dégénéré.
\end{questions}
\label{figlivre-r04-14-var-cor}
\figureLivre[corrige]{r04-14-var}

\label{figlivre-r04-14-cylindre-cor}
\figureLivre[corrige]{r04-14-cylindre}

\end{corrige}
\end{blocLivre}
\begin{blocLivre}{cor-4-15}
\begin{corrige}{15}{Distance minimale à un point}
\demarche{Minimiser d’abord la distance au carré, plus simple à dériver.}
\begin{questions}
\item $AM^2=(x-0)^2+(x^2-2)^2=x^2+x^4-4x^2+4
=\res{x^4-3x^2+4}$.
\item La racine carrée est strictement croissante : les valeurs de $AM=\sqrt{D(x)}$
sont minimales exactement lorsque $D$ l'est.
\item $D'(x)=4x^3-6x=2x(2x^2-3)$.
En posant $a=\sqrt{3/2}$, les signes successifs sont $-$, $+$, $-$, $+$
sur les intervalles séparés par $-a$, 0 et $a$.
$D$ décroît de 8 à $7/4$, croît jusqu'à 4, décroît à $7/4$, puis croît jusqu'à 8.
\item Les points optimaux sont $\res{M_1(-\sqrt{3/2};3/2)}$ et
$\res{M_2(\sqrt{3/2};3/2)}$.
\item La distance minimale vaut $\res{\sqrt7/2\approx1{,}323}$,
strictement inférieure à $AO=2$.
\end{questions}
\label{figlivre-r04-15-var-cor}
\figureLivre[corrige]{r04-15-var}

\label{figlivre-r04-15-distance-cor}
\figureLivre[corrige]{r04-15-distance}

\end{corrige}
\end{blocLivre}
